% ============================================================================
% sec3_structural_model.tex — Structural model of cover-bidder deployment
% ============================================================================

\section{A Framework for Cover-Bidder Deployment}
\label{sec:structural}

We develop a framework for cover bidding in procurement auctions
that formalizes the cover-bidder deployment decision, characterizes
equilibrium cover-bid distributions, and generates testable predictions
that map directly to our empirical strategy
(Section~\ref{sec:empirical}). The framework has three aims: (i)~to
provide economic foundations for the frequent-loser screen; (ii)~to
specify a mixture likelihood that separates cover bids from genuine bids
in the data; and (iii)~to derive the market-selection prediction---that
cover bidding concentrates in competitive markets---as an endogenous
outcome of the cartel's optimization.
Together, these enable \emph{regime identification}---determiningwhether observed FL bids follow the complementary (Regime~1) or
coordinated (Regime~2) cover-bidding equilibrium---which is the
structural model's primary empirical contribution. Unlike
\citet{bajari2003deciding}, who test for collusion without specifying
the cover-bid distribution, our framework uses the FL classification as
a known partition, eliminating the need for EM estimation and directly
identifying the equilibrium cover-bid regime.

% ── 3.1 Environment ──────────────────────────────────────────────

\subsection{Environment}
\label{sec:environment}

Consider a procurement auction for item $i$ at purchasing unit $k$ in
year $t$. The procuring agency posts a reference price $r > 0$ and
selects the lowest bidder.%
\footnote{Under BEC's procurement rules, the reference price is
set by the purchasing unit based on market research---typically
the median of three price quotations from suppliers or the
average of prices in recent comparable tenders. The reference
price is public information, posted before bidding opens. This
administrative rule creates an anchor for both genuine and
cover bidders: the cartel can calibrate $b^*$ just below $r$,
and cover bidders bid above $b^*$ but within a plausible range
of $r$.} There are $n \geq 1$ \emph{genuine} bidders,
each with private cost $c_j \sim F(\cdot \mid \mathbf{x}_j)$, where
$\mathbf{x}_j$ is a vector of observable cost shifters (firm size, CNAE
sector, geographic distance to PBU~$k$). Genuine bidders submit bids
$b_j = c_j + \mu_j$, where $\mu_j \geq 0$ is a markup that depends on
the competitive environment.

A cartel controls firm $j^*$---the designated winner---and deploys
$m \geq 0$ \emph{cover bidders} (frequent losers). Each cover bidder
submits a bid $b_\ell$ drawn from a distribution $G(\cdot \mid b^*,
\boldsymbol{\theta})$, where $b^*$ is the cartel's designated winning
bid and $\boldsymbol{\theta}$ is a vector of strategic parameters. Cover
bids are designed to \emph{not} win: $b_\ell > b^*$ with probability
one. The procuring agency observes $n + m$ total bids.

Two institutional features of S\~ao Paulo's BEC platform shape the
model:

\begin{enumerate}
  \item \textbf{Convite} (sealed-bid): requires at least three bidders
    (Lei~8.666/93, Art.~22, \S7). If $n < 3$, the cartel must deploy
    $m \geq 3 - n$ cover bidders to satisfy the minimum-bidder
    constraint. Bids are submitted simultaneously and opened together.
  \item \textbf{Preg\~ao} (electronic reverse auction): no
    minimum-bidder requirement. Bidders submit descending bids in real
    time. Cover bidders can observe the auction dynamics before
    calibrating their bids, facilitating coordination.
\end{enumerate}

\noindent We denote the auction modality by $\tau \in \{\text{convite},
\text{preg\~ao}\}$, the minimum-bidder constraint by $\underline{n}(\tau)$
(equal to 3 for convite, 0 for preg\~ao), and the detection probability
by $\theta_k \in [0,1]$, which varies across PBUs and is increasing in
oversight capacity.

% ── 3.2 Cartel's Problem ─────────────────────────────────────────

\subsection{The Cartel's Optimization Problem}
\label{sec:cartel_problem}

The cartel chooses the number of cover bidders $m$, the winning bid
$b^*$, and the cover-bid distribution $G$ to maximize expected profit
net of cover-bidding costs and detection penalties:
\begin{equation}
  \max_{m \in \mathbb{Z}_+,\, b^*,\, G} \; \underbrace{\pi(b^*, m, n)
  }_{\text{cartel surplus}}
  - \underbrace{C(m)}_{\text{cover cost}}
  - \underbrace{\theta_k \cdot \Phi(m)}_{\text{expected penalty}}
  \label{eq:cartel_objective}
\end{equation}
subject to:
\begin{align}
  b^* &\leq r \label{eq:ref_price_constraint} \\
  \Pr(b_\ell > b^* \mid G) &= 1 \quad \forall \, \ell \in \{1, \ldots, m\}
    \label{eq:cover_above} \\
  n + m &\geq \underline{n}(\tau) \label{eq:min_bidder}
\end{align}

\noindent We maintain the following assumptions throughout:

\begin{enumerate}
  \item[\textbf{A1.}] Cartel surplus $\pi(b^*, m, n) = b^* - c_{j^*}$ is weakly increasing in $m$ ($\partial \pi / \partial m \geq 0$) with diminishing marginal returns.
  \item[\textbf{A2.}] Diminishing returns to cover bidding: $\partial^2 \pi / \partial m^2 < 0$.
  \item[\textbf{A3.}] Linear cover-bidding cost: $C(m) = c_0 + c_1 m$ with $c_0 \geq 0$ and $c_1 > 0$.
  \item[\textbf{A4.}] Detection penalty $\theta_k \cdot \Phi(m)$ with $\Phi'(m) > 0$ and $\Phi''(m) \geq 0$; for closed-form results we specialize to $\Phi(m) = \phi_0 m$, $\phi_0 > 0$.%
    \footnote{See Appendix~\ref{app:proofs} for the more general specification $\Phi(m, b^*) = \phi_0 m + \phi_1 (b^* - c_{j^*})$, which nests markup-dependent penalties and preserves the comparative statics.}
  \item[\textbf{A5.}] Interior solution exists: $\partial \pi / \partial m \big|_{m=0} > c_1 + \theta_k \phi_0$.
\end{enumerate}

\begin{proposition}[Optimal number of cover bidders]
\label{prop:optimal_m}
Under Assumptions A1--A5, the cartel's optimal cover-bidder count is:
\begin{equation}
  m^* = \max\!\left\{\underline{n}(\tau) - n, \;
  \left\lfloor \frac{\partial \pi / \partial m}{c_1 + \theta_k \phi_0}
  \right\rfloor \right\}
  \label{eq:optimal_m}
\end{equation}
\noindent where the first argument is the corner solution from the
minimum-bidder constraint~\eqref{eq:min_bidder} and the second is the
interior optimum. The optimal $m^*$ satisfies:
\begin{enumerate}
  \item[(i)] $\partial m^* / \partial \theta_k < 0$\,: higher detection
    probability reduces cover bidding;
  \item[(ii)] $\partial m^* / \partial c_1 < 0$\,: higher per-bidder
    cost reduces cover bidding;
  \item[(iii)] $m^*$ is weakly larger under convite than preg\~ao,
    holding $n$ fixed, because $\underline{n}(\text{convite}) = 3 >
    0 = \underline{n}(\text{preg\~ao})$.
\end{enumerate}
\end{proposition}

\begin{proof}
See Appendix~\ref{app:proofs}.
\end{proof}

% ── 3.3 Cover-Bid Distributions ──────────────────────────────────

\subsection{Equilibrium Cover-Bid Distributions}
\label{sec:cover_distributions}

The cartel chooses the cover-bid distribution $G$ to minimize the
probability of detection while satisfying
constraint~\eqref{eq:cover_above}. We characterize two equilibrium
regimes that differ in the cartel's information and coordination
capacity.

\begin{proposition}[Regime~1: Complementary cover bidding]
\label{prop:regime1}
Suppose the cartel can instruct cover bidders to ``bid above $b^*$''
with support $[b^*, b^* + \delta]$ for some $\delta > 0$, but cannot
specify exact bid levels (limited coordination). Suppose further that
the detection algorithm is a variance screen that flags tenders where
the bid distribution departs from maximum entropy over the feasible
support.\footnote{A variance screen computes the Kullback--Leibler
divergence $D_{\text{KL}}(G \| U)$ between the submitted bid
distribution $G$ and the uniform distribution $U$ over $[b^*, b^* +
\delta]$. The cartel minimizes $D_{\text{KL}}$ subject to
$\text{supp}(G) \subseteq [b^*, b^* + \delta]$; the solution is
$G = U$.} Then the equilibrium cover-bid distribution is:
\begin{equation}
  G_1(b) = \frac{b - b^*}{\delta}, \quad b \in [b^*, b^* + \delta]
  \label{eq:regime1}
\end{equation}
Under Regime~1, $\text{Var}(b_\ell) = \delta^2 / 12$. If $\delta$
is large, $\sigma_{\text{cover}} > \sigma_{\text{genuine}}$.
\end{proposition}

\begin{proof}
See Appendix~\ref{app:proofs}.
\end{proof}

\begin{proposition}[Regime~2: Coordinated cover bidding]
\label{prop:regime2}
Suppose the cartel observes $b^*$ in real time (feasible under
preg\~ao) and can specify cover bids precisely (full coordination).
Suppose the detection algorithm is a dispersion screen that flags
tenders where the within-tender bid standard deviation is anomalous.
Then the equilibrium cover-bid distribution solves:
\begin{equation}
  \min_{\sigma_c > 0} \; \left| \sigma_c - \sigma_{\text{genuine}}
  \right| \quad \text{s.t.} \quad \Pr(b_\ell > b^*) = 1
  \label{eq:regime2_opt}
\end{equation}
The solution is a truncated normal distribution centered at
$\mu_c = b^* + \epsilon$ with $\epsilon > 0$ and
$\sigma_c < \sigma_{\text{genuine}}$:
\begin{equation}
  G_2(b) = \Phi\!\left(\frac{b - \mu_c}{\sigma_c}\right),
  \quad \mu_c = b^* + \epsilon, \quad
  \sigma_c < \sigma_{\text{genuine}}
  \label{eq:regime2}
\end{equation}
\end{proposition}

\noindent Under Regime~2, cover bids exhibit lower dispersion than genuine bids, reversing the conventional association between collusion and high bid dispersion.

\begin{proof}
See Appendix~\ref{app:proofs}.
\end{proof}

% ── 3.4 Market Selection ─────────────────────────────────────────

\subsection{Market Selection: Why Cover Bidding Concentrates in
Competitive Markets}
\label{sec:market_selection}

\begin{proposition}[Market-selection prediction]
\label{prop:market_selection}
Under Assumptions A1--A2, if the marginal return to cover bidding is
decreasing in market concentration---formally, if
\begin{equation}
  \frac{\partial^2 \pi}{\partial m \, \partial \text{HHI}} < 0
  \label{eq:cross_deriv}
\end{equation}
where $\text{HHI}$ is the Herfindahl--Hirschman Index of genuine
competition---then $m^*$ is decreasing in HHI. Cover bidding
concentrates in competitive markets.
\end{proposition}

\noindent Cover bidding thus acts as a strategic complement to genuine competition.

\begin{proof}
See Appendix~\ref{app:proofs}.
\end{proof}

% ── 3.5 Structural Likelihood ────────────────────────────────────

\subsection{Structural Likelihood}
\label{sec:likelihood}

Each bid $b_{jt}$ is drawn from one of two distributions, with the FL classification providing known group membership (no EM required):
\begin{equation}
  b_{jt} \sim
  \begin{cases}
    F_{\text{genuine}}(b \mid \mathbf{x}_j, \boldsymbol{\alpha}) &
      \text{if firm } j \text{ is genuine (non-FL)} \\
    G_{\text{cover}}(b \mid b^*_t, \boldsymbol{\gamma}) &
      \text{if firm } j \text{ is FL}
  \end{cases}
  \label{eq:mixture}
\end{equation}
Genuine bids are log-normal with item and year fixed effects:
\begin{equation}
  \log b_{jt} \mid j \text{ genuine} \sim
  N\!\left(\mathbf{x}_j'\boldsymbol{\beta} + \alpha_i + \lambda_t,
  \; \sigma_g^2\right)
  \label{eq:genuine_dist}
\end{equation}
Cover bids follow $\text{Uniform}[b^*_t, b^*_t + \delta]$ under Regime~1:
\begin{equation}
  b_{jt} \mid j \text{ FL} \sim
  \text{Uniform}[b^*_t, \; b^*_t + \delta]
  \label{eq:cover_dist_r1}
\end{equation}
and a truncated normal under Regime~2:
\begin{equation}
  \log b_{jt} \mid j \text{ FL} \sim
  N\!\left(\log b^*_t + \epsilon, \; \sigma_c^2\right),
  \quad \sigma_c < \sigma_g
  \label{eq:cover_dist_r2}
\end{equation}
The tender-level log-likelihood is:
\begin{equation}
  \ell_t = \sum_{j \in \mathcal{G}_t} \log f_{\text{genuine}}(b_{jt}
  \mid \mathbf{x}_j, \boldsymbol{\alpha})
  + \sum_{\ell \in \mathcal{F}_t} \log g_{\text{cover}}(b_{\ell t}
  \mid b^*_t, \boldsymbol{\gamma})
  \label{eq:loglik}
\end{equation}

% ── 3.6 Identification ───────────────────────────────────────────

\subsection{Identification of Structural Parameters}
\label{sec:identification}

The genuine-bid parameters $(\boldsymbol{\beta}, \sigma_g^2)$ are identified from within-item--year variation in non-FL losing bids. The cover-bid parameters $(\delta$ or $\sigma_c)$ are identified from the distribution of FL bids relative to the winning bid $b^*_t$ in each tender. The detection probability $\theta_k$ is partially identified from cross-PBU variation in oversight capacity, proxied by procurement volume. The cartel markup itself is not identified by the distributional parameters but by the reduced-form comparison of FL-present and FL-absent tenders (Section~\ref{sec:results_structural}).

% ── 3.7 Testable Predictions ─────────────────────────────────────

\subsection{Testable Predictions}
\label{sec:predictions}

The model generates five testable predictions, each mapped to a specific
empirical test in Section~\ref{sec:empirical}:

\begin{prediction}[Price effect]
\label{pred:price_v6}
FL-present tenders have higher winning prices than comparable tenders
without FL, controlling for item type, year, and PBU.
\emph{Test:} OLS regression (Section~\ref{sec:reduced_form}).
\end{prediction}

\begin{prediction}[Strategic selection]
\label{pred:selection_v6}
FL-present tenders have at least as many genuine (non-FL) competitors as
FL-absent tenders. Cover bidders are deployed on top of genuine
competition, not as substitutes.
\emph{Test:} Non-FL firm count regression (Section~\ref{sec:reduced_form}).
\end{prediction}

\begin{prediction}[Regime identification]
\label{pred:regime_v6}
Under Regime~1, FL bid dispersion exceeds genuine bid dispersion
($\sigma_{\text{cover}} > \sigma_{\text{genuine}}$). Under Regime~2,
FL bid dispersion is lower ($\sigma_{\text{cover}} <
\sigma_{\text{genuine}}$). BIC model selection determines the
empirically dominant regime.
\emph{Test:} Structural likelihood comparison and bid-dispersion
regression (Section~\ref{sec:structural_estimation}).
\end{prediction}

\begin{prediction}[Market selection]
\label{pred:market_v6}
The FL--price effect is decreasing in market concentration (HHI).
Cover bidding concentrates in competitive markets where the marginal
return to manufactured competition is highest.
\emph{Test:} Network-split regression and FL $\times$ HHI interaction
(Section~\ref{sec:empirical_network}).
\end{prediction}

\begin{prediction}[Exchangeability and conditional-independence violation]
\label{pred:exchange_v6}
FL bid residuals are non-exchangeable with non-FL bid residuals
conditional on observables, and FL residuals within the same tender
exhibit significantly positive pairwise correlation. The structural
model implies that FL bids are drawn from $G_{\text{cover}}$ rather
than $F_{\text{genuine}}$, producing distinguishable residual
distributions and within-tender dependence. Under tender FE, which
absorb the shared focal-point component inflated by cover bidding,
the relative magnitude of FL versus non-FL correlation is ambiguous.
\emph{Tests:} Kolmogorov--Smirnov test and pairwise-product test on
Bajari--Ye first-stage residuals (Section~\ref{sec:bajari_ye}).
\end{prediction}
